%%%%%%%%%%%%%%%%%%%%%%%%%%%%%%%%%%%%%%%%%%%%%%%%%%%%%%%%%%%%%
%                                                           %
% BEGINNING DATE: 22 - October - 1998                       %
% LAST MODIFICATION DATE: 11 - abril - 2000                 %
%                                                           %
% Written with LaTeX2e                                      %
% AMS package version 2.0                                   %
%                                                           %
%%%%%%%%%%%%%%%%%%%%%%%%%%%%%%%%%%%%%%%%%%%%%%%%%%%%%%%%%%%%%

\documentclass[12pt,leqno]{amsart}
\textwidth=30cc
\baselineskip=16pt

\theoremstyle{plain}
\newtheorem{Thm}{Theorem}[section]
\newtheorem{Cor}[Thm]{Corollary} % Same numbering sequence that Thm
\newtheorem{Pro}[Thm]{Proposition}
\newtheorem{Lem}[Thm]{Lemma}

\newcommand{\supp}{\operatorname{supp}} % support
\newcommand{\lpnorm}[3]{\|{#1}\|_{L^{#2}({#3})}} % L^p-norm of a function

\newcommand{\Id}{\operatorname{Id}} % identity map
\newcommand{\spa}{\operatorname{span}} % span
\newcommand{\htr}{\mathcal{H}} % Hankel transform
\newcommand{\htra}{\mathcal{H}_\alpha} % Hankel transform of order \alpha
\newcommand{\hab}{\mathcal{H}_{\alpha,\beta}}
\newcommand{\hf}{\mathop{{}_2F_1}} % hypergeometric function
\newcommand{\ab}{{(\alpha,\beta)}} % An alias

\date{PUBLISHED IN \textit{Studia Math.}
\textbf{142} (2000), no.~2, 253--267.}

%% A fix (suggested by tech-support@ams.org) to avoid that
%% an overfull appears in the heading of every even page
%% when we have an accent in the name of an author.
%% (Put it before \begin{document})
%\makeatletter
%\renewcommand{\leftmark}{%
%  \vbox to\headheight{\vss
%    \hbox{\expandafter \@firstoftwo \topmark {}{}}%
%  }%
%}
%\makeatother

%%%%%%%%%%%%%%%%%%%%%%%%%%%%%%%%%%%%%%%%%%%%%%%%%%%%%%%%%%%%%%%%%%%%%%%

\begin{document}

\title[Solving Dual Integral Equations on Lebesgue Spaces]{Solving Dual
Integral Equations on Lebesgue Spaces}

\author[\'O. Ciaurri]{\'Oscar Ciaurri}
\address{Departamento de Matem\'aticas y Computaci\'on,
Universidad de La Rioja, Edificio J.~L.~Vives,
Calle Luis de Ulloa s/n, 26004 Logro\~no, Spain}
\email{osciaurr@dmc.unirioja.es}
%\thanks{Research of the first author supported by a grant of the UR}

\author[J. J. Guadalupe]{Jos\'e J. Guadalupe}
\address{Departamento de Matem\'aticas y Computaci\'on,
Universidad de La Rioja, Edificio J.~L.~Vives,
Calle Luis de Ulloa s/n, 26004 Logro\~no, Spain}
\email{joguadal@dmc.unirioja.es}

\author[M. P\'erez]{Mario P\'erez}
\address{Departamento de Matem\'aticas,
Universidad de Zaragoza, Edificio de Matem\'aticas,
Ciudad Universitaria s/n, 50009 Zaragoza, Spain}
\email{mperez@posta.unizar.es}

\author[J. L. Varona]{Juan~L.~Varona}
\address{Departamento de Matem\'aticas y Computaci\'on,
Universidad de La Rioja, Edificio J.~L.~Vives,
Calle Luis de Ulloa s/n, 26004 Logro\~no, Spain}
\email{jvarona@dmc.unirioja.es}
\urladdr{http://www.unirioja.es/dptos/dmc/jvarona/welcome.html}
\thanks{Research supported by grants of the DGES and UR}

\thanks{On April 1, 2000, while this paper was being revised, Jos\'e
J.~Guadalupe unexpectedly died. We still can not say how deep is our sorrow}

\subjclass[2000]{Primary 45F10; Secondary 42C10}

\keywords{Dual integral equations, Bessel functions, Fourier series, 
Hankel transform.}

\begin{abstract}
We study dual integral equations associated with Hankel transforms, that is,
dual integral equations of Titchmarsh's type.  We reformulate these equations
giving a better description in terms of continuous operators on $L^p$ spaces,
and we solve them in these spaces. The solution is given both as an operator
described in terms of integrals and as a series $\sum_{n=0}^\infty c_n
J_{\mu+2n+1}$ which converges in the $L^p$-norm and almost everywhere, where
$J_\nu$ denotes the Bessel function of order $\nu$. Finally, we study the
uniqueness of the solution.
\end{abstract}

%%%%%%%%%%%%%%%%%%%%%%%%%%%%%%%%%%%%%%%%%%%%%%%%%%%%%%%%%%%%%%%%%%%%%%%
\maketitle
%%%%%%%%%%%%%%%%%%%%%%%%%%%%%%%%%%%%%%%%%%%%%%%%%%%%%%%%%%%%%%%%%%%%%%%

%%%%%%%%%%%%%%%%%%%%%%%%%%%%%%%%%%%%%%%%%%%%%%%%%%%%%%%%%%%%%%%%%%%%%%%
\section{Introduction}
%%%%%%%%%%%%%%%%%%%%%%%%%%%%%%%%%%%%%%%%%%%%%%%%%%%%%%%%%%%%%%%%%%%%%%%

In some physical problems related with potential and electromagnetic or
acoustic radiation theory, sometimes the unknown function satisfies an
integral equation over part of the range $(0,\infty)$ and a different integral
equation over the rest of the range. These equations are known as
dual integral equations. An important case is the so-called dual
integral equations of Titchmarsh's type:
\begin{equation}
    \label{ec:DIEorig}
    \left\{
      \begin{array}{ll}
      \displaystyle
      \int_0^\infty t^\beta f(t) J_\alpha(xt)\,dt = g(x)&
        \text{if $0<x<1$,} \\[8pt]
      \displaystyle
      \int_0^\infty f(t)J_\alpha(xt)\,dt = 0&
        \text{if $x>1$,}
      \end{array}
    \right.
\end{equation}
where $J_\alpha$ stands for the Bessel function of order $\alpha$ (see
\cite{Wat} or \cite[Ch.~VII]{E-HTF}), $g$ is a given function and $f$ is the
unknown function. For a function~$h$,
\begin{equation*}
    \int_0^\infty h(t)J_\alpha(xt) (xt)^{1/2} \,dt, \quad x > 0,
\end{equation*}
is usually known as the Hankel transform of~$h$; so, the second equation
in~(\ref{ec:DIEorig}) means that the Hankel transform of~$t^{-1/2} f(t)$ is
supported on $[0,1]$, and the first one imposes a condition on the Hankel
transform of~$t^{\beta-1/2} f(t)$.

There are different methods to solve these equations, most of them only formal. 
For instance, they can be solved by using Mellin transforms or some other
integral transforms. Also, they can be reduced to Fredholm integral equations.
Usually, these methods allow to find the solution $f$ as an explicit
expression with integrals; we can see some of them in the
books~\cite[p.~337]{Tit}, \cite[\S~12, p.~65]{Sne}, \cite[\S~5.11]{Hoc}
and~\cite[p.~76]{E-HTF}. Another method consists of solving the equation as a
series $\sum_{n=0}^\infty c_n J_{\mu+2n+1}$; see~\cite{Tra1} and~\cite{Tra2},
the first one with a large bibliography. But, as long as the authors know, it
is only studied as a formal method.

In this paper we pursuit a rigorous approach to solve dual integral equations.
We reformulate~(\ref{ec:DIEorig}) so as to obtain a better description in
terms of operators on $L^p$ spaces, and we find the solution in these
spaces. Also, we identify the solution as a Fourier-Neumann series whose $L^p$
and almost everywhere convergence is studied.

The paper is organised as follows: in section~\ref{sec:die} we state the dual
integral equation in a more convenient form and define some associated
operators. Section~\ref{sec:BfJp} collects some properties of Bessel functions
and Jacobi polynomials. We describe a solution to the dual integral equation
in section~\ref{sec:solution}, and section~\ref{sec:uniq} is devoted to the
uniqueness of the solution. Sections~\ref{sec:proofHTjn}
and~\ref{sec:proofLab} contain some of the proofs.

Throughout this paper, unless otherwise stated, we will use $C$ to denote a
positive constant independent of $f$ (and all other variables), which can
assume different values in different occurrences.

Also, for any function $g$ defined on $[0,1]$, the extension given by $g(x) =
0$ at each $x > 1$ will be denoted by $\chi_{[0,1]} g$, with a small abuse of
notation. Strictly speaking, $\chi_{[0,1]}$ could be understood either as a
characteristic function or as an operator taking functions defined on $[0,1]$
to functions defined on $[0,\infty)$.

%%%%%%%%%%%%%%%%%%%%%%%%%%%%%%%%%%%%%%%%%%%%%%%%%%%%%%%%%%%%%%%%%%%%%%%
\section{The dual integral equation}
\label{sec:die}
%%%%%%%%%%%%%%%%%%%%%%%%%%%%%%%%%%%%%%%%%%%%%%%%%%%%%%%%%%%%%%%%%%%%%%%

Let us define, for $\alpha > -1$, the integral operator $\htra$ as
\begin{equation*}
    \htra(f,x) = \frac{x^{-\alpha/2}}{2}
     \int_0^\infty {f(t)J_\alpha (\sqrt {xt})t^{\alpha/2}\,dt},
     \quad x > 0,
\end{equation*}
for suitably integrable functions $f$. For instance, $\htra$ is an
isomorphism from the Schwartz class
\begin{equation*}
    S^+ = \{ f \in C^\infty((0,\infty)) : \forall k,j \ge 0,
    \ |t^kf^{(j)}(t)| < C_{k,j} \}
\end{equation*}
onto itself and $\htra^2$ is the identity map. This operator is a modified
Hankel transform. For $\alpha \geq -1/2$, $1 \leq p \leq 2$, and $1/p + 1/p'
= 1$, $\htra$ extends to a bounded operator from
$L^p([0,\infty),x^\alpha\,dx)$ into $L^{p'}([0,\infty),x^\alpha\,dx)$, i.e.
\begin{equation*}
    \lpnorm{\htra f}{p'}{[0,\infty),x^\alpha dx}
      \le C \lpnorm{f}p{[0,\infty),x^\alpha dx},
      \quad f \in L^p([0,\infty),x^\alpha\,dx).
\end{equation*}
However, the Hankel transform does not extend to
$L^p([0,\infty),x^\alpha\,dx)$ if $2 < p$ (see~\cite{Dur, St-Tr, Var}).

Another operator will be used: the multiplier of the Hankel transform
associated to $\chi_{[0,1]}$, that is, the operator $M_\alpha$ given by
$\htra(M_\alpha f) =  \chi_{[0,1]} \htra f$. This multiplier plays an
important role in the study of orthogonal Fourier expansions (see~\cite{Var}
in connection with Fourier-Neumann series, and~\cite{Ste} for Laguerre
series).

Herz's classical result determines the range of $p$ such that $M_\alpha$ is a
well defined, bounded operator from $L^p([0,\infty), x^\alpha \, dx)$ into
itself (\cite{Her}; see also~\cite{Ste, Var}):

\begin{Pro}
\label{pro:acotMa}
Let $\alpha \ge -1/2$ and $1<p<\infty$. Then
\begin{equation*}
    \lpnorm{M_{\alpha}f}p{[0,\infty),x^{\alpha}dx}
      \le C \lpnorm{f}p{[0,\infty),x^{\alpha}dx} 
    \iff
    \tfrac{4(\alpha+1)}{2\alpha+3} < p < \tfrac{4(\alpha+1)}{2\alpha+1}.
\end{equation*}
\end{Pro}

For more general results on Hankel multipliers, see~\cite{St-Tr} and the
references therein.

In a dense subset of $L^p([0,\infty),
x^\alpha
\, dx)$ (for instance, $S^+$)
\begin{equation*}
    M_{\alpha}f = \htra(\chi_{[0,1]}\htra f)
\end{equation*}
and $\htra^2 = \Id$. Whenever $\htra$ is well defined, it follows that $\htra
f$ is supported on $[0,1]$ if and only if $M_\alpha f = f$.

Now, let us reformulate the dual integral equations.
With a simple change of notation, we can write~(\ref{ec:DIEorig}) as
\begin{equation}
    \label{ec:DIEmodi}
    \left\{
      \begin{array}{ll}
      \displaystyle \frac{x^{-\alpha/2}}{2} \int_0^\infty t^{\beta} f(t)
        J_\alpha(\sqrt{xt}) t^{\alpha/2}\,dt = g(x)&
          \text{if $0<x<1$,} \\[9pt]
      \displaystyle \frac{x^{-\alpha/2}}{2} \int_0^\infty f(t)
        J_\alpha(\sqrt{xt}) t^{\alpha/2}\,dt = 0&
          \text{if $x>1$}.
      \end{array}
    \right.
\end{equation}
The second equation in (\ref{ec:DIEmodi}) means that $\supp (\htra f)
\subseteq [0,1]$; in other words, $M_{\alpha}f = f$, provided that $f$
belongs to a convenient $L^p$ space.

The first equation in~(\ref{ec:DIEmodi}) can be read as $\htra(t^\beta f)
\chi_{[0,1]} = \chi_{[0,1]}g$. Under certain conditions, $\htra$ is an
inversible operator. Then, we obtain the equivalent equation
$M_{\alpha}(t^{\beta}f,x) = \htra(\chi_{[0,1]}g,x)$. It will be convenient to
multiply both sides by $x^{-\beta}$, so we get
$x^{-\beta}M_{\alpha}(t^{\beta}f,x) = x^{-\beta}\htra(\chi_{[0,1]}g,x)$.

To sum up, we are interested in solving in $L^p([0,\infty),x^{\alpha}\,dx)$
the equation
\begin{equation}
    \label{ec:DIEoper}
    \left\{
      \begin{array}{l}
         x^{-\beta}M_{\alpha}(t^{\beta}f,x)
           = x^{-\beta}\htra(\chi_{[0,1]}g,x), \\[3pt]
         M_{\alpha}f = f. 
      \end{array}
    \right.
\end{equation}

In a strict sense,~(\ref{ec:DIEmodi}) and~(\ref{ec:DIEoper}),
are not exactly equivalent if we do not assume that the functions belong to a
convenient $L^p$ space. However, it is interesting to note that, for any
practical physical application, the interpretation of a dual integral
equation and its solution as in~(\ref{ec:DIEmodi}) is equivalent to its
interpretation as in~(\ref{ec:DIEoper}).

Together with $\htra$ and $M_{\alpha}f = \htra(\chi_{[0,1]}\htra f)$, let us
consider the operators $M_{\alpha,\beta}$ and $\hab$ given by
\begin{align*}
    M_{\alpha,\beta}f &= x^{-\beta}M_{\alpha}(t^{\beta}f),\\
    \hab g &= x^{-\beta}\htra(\chi_{[0,1]}g).
\end{align*}
With this notation, the dual integral equation~(\ref{ec:DIEoper}) can be
written as
\begin{equation}
    \label{ec:DIEfinal}
    \left\{
      \begin{array}{l}
         M_{\alpha,\beta}f = \hab g, \\[3pt]
         M_{\alpha}f = f.
      \end{array}
    \right.
\end{equation}
Those operators are well defined, for instance, if $f \in S^+$ and $g \in
C^\infty([0,1])$. We see below that $M_{\alpha,\beta}$ is bounded with the
$L^p([0,\infty), x^\alpha \, dx)$-norm, under some conditions on $\alpha$,
$\beta$, and $p$. Therefore, it extends to a bounded operator on
$L^p([0,\infty), x^\alpha \, dx)$. With a similar argument, $\hab$ extends to
a bounded operator from $L^p([0,1], x^\alpha \, dx)$ into $L^p([0,\infty),
x^\alpha \, dx)$.

\begin{Pro}
\label{pro:acotMab}
Let $\alpha \ge -1/2$, $\beta \ge 0$ and $1<p<\infty$. Then
\begin{equation*}
    \lpnorm{M_{\alpha,\beta}f}p{[0,\infty),x^{\alpha}dx}
      \le C \lpnorm{f}p{[0,\infty),x^{\alpha}dx} 
    \iff
    \tfrac{4(\alpha+1)}{2\alpha+4\beta+3} < p
      < \tfrac{4(\alpha+1)}{2\alpha+4\beta+1}.
\end{equation*}
\end{Pro}
\begin{proof}
We give only a sketch of the proof. It follows the proof for $M_\alpha$
in~\cite{Var}. Actually, this is a particular case of weighted versions of
Herz's classical result (with power weights).

Let us take $p_0 = \tfrac{4(\alpha+1)}{2\alpha+4\beta+3}$ and
$p_1 = \tfrac{4(\alpha+1)}{2\alpha+4\beta+1}$. For each $f \in S^+$, Fubini's
theorem applies to $M_{\alpha,\beta}f =
x^{-\beta}\htra(\chi_{[0,1]}\htra(t^{\beta}f))$, then Lommel's
formula
\begin{equation*}
    \int_0^1 J_\alpha(yt)J_\alpha(yx) y\,dy
    = \frac1{t^2-x^2} (tJ_{\alpha+1}(t)J_{\alpha}(x)
      - xJ_{\alpha}(t)J_{\alpha+1}(x))
\end{equation*}
gives
\begin{align*}
   M_{\alpha,\beta}(f,x) &=
     \tfrac12 x^{-\alpha/2-\beta+1/2} J_{\alpha+1}(x^{1/2})
     H(t^{\alpha/2+\beta}J_{\alpha}(t^{1/2})f(t),x)
    \\
        &- \tfrac12 x^{-\alpha/2-\beta} J_{\alpha}(x^{1/2})
     H(t^{\alpha/2+\beta+1/2}J_{\alpha+1}(t^{1/2})f(t),x)
   \\
    &= W_1(f,x) - W_2(f,x),
\end{align*}
where $H$ is the Hilbert transform $H(f,x) = \int_0^\infty
\frac{f(t)}{x-t}\,dt$. The Hilbert transform is a bounded operator from
$L^p([0,\infty),x^\lambda\,dx)$ into itself if and only if $-1<\lambda<p-1$.
Fix $1 < p < \infty$; then, by using the bound $|J_{\alpha}(x)| \le C
x^{-1/2}$, it is easy to check that $W_1$ and $W_2$ are bounded operators on
$L^p([0,\infty),x^\alpha\,dx)$ if
\begin{equation*}
   p_0 < p < \tfrac{4(\alpha+1)}{2\alpha+4\beta-1}
\end{equation*}
(disregard the right hand side inequality if $2\alpha+4\beta-1 \leq 0$)
and
\begin{equation*}
   \tfrac{4(\alpha+1)}{2\alpha+4\beta+5} < p < p_1,
\end{equation*}
respectively. Then, $M_{\alpha,\beta}$ is bounded if $p_0 < p < p_1$.

In fact, $p_0<p<p_1$ is a necessary condition for the boundedness of
$M_{\alpha,\beta}$. By interpolation, we only need to observe that
$M_{\alpha,\beta}$ is not bounded for $p=p_0$ (if $p_0 > 1$) and $p=p_1$. If
$p = p_0 > 1$, $W_2$ is bounded; however, more precise estimates for the
Bessel functions near infinity and a clever election of $f$ show that $W_1$
is not bounded. Then, $M_{\alpha,\beta}$ is not bounded. The case of $p =
p_1$ is analogous.
\end{proof}

Regarding the Hankel transform $\htra$, we have the following theorem of
Rooney (\cite[p.~1100]{Roo},~\cite{H-R}, after a change of notation):

\begin{Thm}[Rooney]
\label{th:Rooney}
Let $\alpha>-1$, $1 < p \leq q < \infty$, $\max\{\tfrac1p,1-\tfrac1{q}\} \le
\nu < \alpha+\tfrac32$. Then
\begin{equation*}
    \Big( \int_0^{\infty} |x^{-\nu/2+\alpha/2+3/4}\htra(h,x)|^q
        \, \frac{dx}x \Big)^{1/q}
      \le C \Big( \int_0^{\infty} |x^{\nu/2+\alpha/2+1/4}h(x)|^p \,\frac{dx}x
        \Big)^{1/p}.
\end{equation*}
\end{Thm}

The boundedness of $\hab$ follows as a consequence:

\begin{Pro}
\label{pro:acotHab}
Let $\alpha \ge -1/2$, $\beta \ge 0$, $1<p<\infty$ and assume
\begin{equation*}
    \frac{2(\alpha+1)}{\alpha+\beta+1} \le p < \frac{\alpha+1}{\beta}.
\end{equation*}
Then $\lpnorm{\hab g}p{[0,\infty),x^{\alpha}dx} \le C
\lpnorm{g}p{[0,1],x^{\alpha}dx}$.
\end{Pro}

\begin{proof} Take $\nu = 2 \beta + \alpha + \tfrac32 - \tfrac{2(\alpha+1)}p$
and $p=q$. It is easy to see that we can apply Theorem~\ref{th:Rooney} and get
\begin{align*}
   \lpnorm{\hab g}p{[0,\infty),x^{\alpha}dx}
   &= \lpnorm{x^{-\beta+(\alpha+1)/p}
   \htra(\chi_{[0,1]}g)}p{[0,\infty),\frac{dx}x}
   \\
   &\le C
   \lpnorm{x^{\beta+(\alpha+1)(1-1/p)}\chi_{[0,1]}g}p{[0,\infty),\frac{dx}x}
   \\
    &= C
  \lpnorm{x^{\beta+(\alpha+1)(1-2/p)}\chi_{[0,1]}g}p{[0,\infty),x^{\alpha}dx}
   \\
    &\le C \lpnorm{g}p{[0,1],x^{\alpha}dx},
\end{align*}
where the last inequality follows from $\beta+(\alpha+1)(1-2/p) \ge 0$.
\end{proof}

Then, our dual integral equation~(\ref{ec:DIEfinal}) is well posed and the
question we try to solve is the following: given any $g \in
L^p([0,1],x^{\alpha}\,dx)$, is there a (unique) solution $f \in
L^p([0,\infty),x^{\alpha}\,dx)$?

%%%%%%%%%%%%%%%%%%%%%%%%%%%%%%%%%%%%%%%%%%%%%%%%%%%%%%%%%%%%%%%%%%%%%%%
\section{Bessel functions and Jacobi polynomials}
\label{sec:BfJp}
%%%%%%%%%%%%%%%%%%%%%%%%%%%%%%%%%%%%%%%%%%%%%%%%%%%%%%%%%%%%%%%%%%%%%%%

If $\alpha>-1$, the Bessel functions satisfy the orthogonality relation
\begin{equation*}
   \int_0^\infty  J_{\alpha +2n+1}(x) J_{\alpha +2m+1}(x)\,\frac{dx}{x}
     = \frac{\delta _{nm}}{2(2n+\alpha +1)},
   \quad n,m=0,1,2,\ldots
\end{equation*}
After a change of variable, the system $\{j_n^\alpha\}_{n=0}^\infty$ given by
\begin{equation*}
    j_n^\alpha(x)
      = \sqrt{\alpha+2n+1} J_{\alpha+2n+1}(\sqrt{x})
        x^{-\alpha/2-1/2}
\end{equation*}
is orthonormal on $L^2([0,\infty),x^\alpha\,dx)$. There is a tight relation
between Bessel functions and Jacobi polynomials $P^\ab_n$ and the following
lemma is relevant for our purposes; the first part was proved in~\cite{CGPV}
and the second part will be proved in section~\ref{sec:proofHTjn}. Of
course, these formulas hold in Lebesgue spaces, that is, almost everywhere.

\begin{Lem}
\label{lem:HTjn}
Let $\alpha, \beta>-1$ with $\alpha+\beta>-1$. Then
\begin{equation}
    \label{ec:HTjnab}
    \htra(j_n^{\alpha+\beta},x)
      = 2^{-\beta} \tfrac{\sqrt{\alpha+\beta+2n+1} \, \Gamma(n+1)}
        {\Gamma(\beta+n+1)}
        (1-x)^\beta P^\ab_n(1-2x) \chi_{[0,1]}(x).
\end{equation}
Assume further $\beta<1$. Then
\begin{equation}
    \label{ec:HTbjnab}
     \chi_{[0,1]}(x) \htra(t^{\beta}j_n^{\alpha+\beta},x)
      = 2^\beta \tfrac{\sqrt{\alpha+\beta+2n+1} \, \Gamma(\alpha+\beta+n+1)}
        {\Gamma(\alpha+n+1)} P^\ab_n(1-2x) \chi_{[0,1]}(x).
\end{equation}
\end{Lem}

In particular, $\supp (\htra(j^{\alpha+\beta}_n)) \subseteq [0,1]$.
However, note that~(\ref{ec:HTbjnab}) refers only to the Hankel transform
of $t^{\beta}j_n^{\alpha+\beta}$ at $x \in [0,1]$; nothing is claimed for
$x>1$.


The Jacobi polynomials $\{P^\ab_n(x)\}_{n=0}^\infty$ of order $\alpha,
\beta$ (see \cite[Ch.~X]{E-HTF} or \cite[Ch.~IV]{Sze}) are orthogonal
on $[-1,1]$ with respect to the weight $(1-x)^\alpha(1+x)^\beta$,
$\alpha,\beta>-1$.

After a change of variable, the system $\{P_n^\ab(1-2x)\}_{n=0}^\infty$ is
orthogonal on $[0,1]$ with respect to the weight $x^\alpha(1-x)^\beta$,
$\alpha,\beta>-1$. To be precise, the orthogonality relation for these
polynomials is
\begin{equation*}
    \int_0^1 P_n^\ab(1-2x) P_m^\ab(1-2x) x^\alpha (1-x)^\beta \,dx
      = h_n^{\ab}\delta_{nm}
\end{equation*}
with 
\begin{equation*}
    h_n^{\ab} = \frac{\Gamma(\alpha+n+1) \Gamma(\beta+n+1)}
        {(\alpha+\beta+2n+1) \Gamma(\alpha+\beta+n+1) n!}\ .
\end{equation*}
Let us take
\begin{equation}
    \label{ec:newja}
    p_n^\ab(x) = (h_n^{\ab})^{-1/2}P_n^{\ab}(1-2x),\quad n=0,1,2,\ldots
\end{equation}
This is an orthonormal system on $[0,1]$ with respect to the weight $w(x)
= x^{\alpha}(1-x)^{\beta}$. For any suitable function~$g$ defined on $[0,1]$,
its Fourier-Jacobi series is the formal expansion
\begin{equation*}
    g \sim \sum_{n=0}^{\infty}a_n(g)p_n^\ab,
    \quad
    a_n(g) = \int_0^1 g(x)p_n^\ab(x)x^{\alpha}(1-x)^{\beta}\,dx.
\end{equation*}

The following result by Muckenhoupt gives conditions for the uniform
boundedness and the mean convergence of $S_n g$ (actually, both are
equivalent, by the Banach-Steinhaus theorem):

\begin{Thm}[Muckenhoupt~\cite{Muc}]
Assume that $\alpha > -1$, $\beta > -1$, $1 < p < \infty$ and let $S_n g$
denote the $n$th partial sum of the Jacobi polynomial series for $g$ with
parameters $\alpha$ and $\beta$. Assume that
\begin{gather*}
   \left| a + (\alpha+1) \Big( {\frac1p-\frac12} \Big) \right|
   < \min\left\{\frac14, \frac{\alpha+1}2\right\},
   \\
   \left| b + (\beta+1) \Big( {\frac1p-\frac12} \Big) \right|
   < \min\left\{\frac14, \frac{\beta+1}2\right\}.
\end{gather*}
Then, there exists a constant $C$ such that
\begin{equation*}
    \lpnorm{x^a (1-x)^b S_ng}p{[0,1],x^\alpha (1-x)^\beta\,dx}
    \leq C \lpnorm{x^a (1-x)^b g}p{[0,1],x^\alpha (1-x)^\beta\,dx}
\end{equation*}
for every $n \in \mathbb{N}$, and
\begin{equation*}
   \lim_n \lpnorm{x^a (1-x)^b (S_ng - g)}p{[0,1],x^\alpha (1-x)^\beta\,dx}
   = 0
\end{equation*}
for every $g$ with $\lpnorm{x^a (1-x)^b g}p{[0,1],x^\alpha (1-x)^\beta\,dx} <
\infty$.
\end{Thm}

For our purposes, we will only need the following:

\begin{Cor}
\label{cor:convja}
Let $\alpha \ge -1/2$, $\beta \ge 0$, $1<p<\infty$ and
\begin{equation*}
    \max{\left\{ \frac{4(\alpha+1)}{2\alpha+3}, \frac4{2\beta+3} \right\}}
    < p <
    \min{\left\{ \frac{4(\alpha+1)}{2\alpha+1}, \frac4{2\beta+1} \right\}}.
\end {equation*}
Then,
\begin{equation*}
    \lim_{n\to \infty}S_ng = g
\end{equation*}
in the $L^p([0,1],x^{\alpha}\,dx)$-norm, for any $g \in
L^p([0,1],x^\alpha\,dx)$.
\end{Cor}
\begin{proof}
Take $a = 0$ and $b = -\beta/p$ in the previous result.
\end{proof}

The scheme we use to solve the dual equation~(\ref{ec:DIEfinal}) goes as
follows: expand $g$ as a Fourier-Jacobi series, that is, $g =
\sum_{n=0}^\infty a_n p_n^\ab$; then, the solution is $f =
\sum_{n=0}^\infty b_n j_n^{\alpha+\beta}$, where $b_n$ is explicitely given
in terms of $a_n$ and the series converges both in $L^p$ and almost
everywhere.

Series of the form $\sum_{n=0}^\infty c_n J_{\mu+n}$ are usually known as
Neumann series. Thus, we are describing the solution of the dual integral
equation as a Fourier-Neumann series.

The operator that takes $g$ into $f$ will be proved to be bounded
on $L^p$. It can also be written in terms of integral operators.

%%%%%%%%%%%%%%%%%%%%%%%%%%%%%%%%%%%%%%%%%%%%%%%%%%%%%%%%%%%%%%%%%%%%%%%%%%%
\section{Main results: The solution of the equation}
\label{sec:solution}
%%%%%%%%%%%%%%%%%%%%%%%%%%%%%%%%%%%%%%%%%%%%%%%%%%%%%%%%%%%%%%%%%%%%%%%%%%%

In this section we introduce the operator $L_{\alpha,\beta}$
and state some of its properties. This operator solves
the dual integral equation~(\ref{ec:DIEfinal}).

Let $g$ a suitable function on $[0,1]$. We define $L_{\alpha,\beta}g$ by
\begin{equation*}
    L_{\alpha,\beta}(g,x) = \frac{1}{2^{\beta+1}\Gamma(\beta)}
    \int_0^1 \frac{J_{\alpha+\beta}(\sqrt{xt})}{(xt)^{{\alpha/2+\beta}/2}}
    \int_0^t g(u)(t-u)^{\beta-1}u^{\alpha}\,du \,dt,
    \quad x>0,
\end{equation*}
if $\beta > 0$ and $L_{\alpha,0}g = \htra(\chi_{[0,1]} g)$. Our first result
states that $L_{\alpha,\beta}$ is a bounded operator from $L^p([0,1],
x^\alpha \, dx)$ into $L^p([0,\infty), x^\alpha \, dx)$:

\begin{Thm}
\label{th:acotLab}
Let $\alpha \ge -1/2$, $\beta \ge 0$, $1 < p < \infty$ and
\begin{equation*}
    \frac{2(2\alpha+3)}{2(\alpha+\beta)+3} \le p < \infty.
\end{equation*}
Then
\begin{equation*}
    \lpnorm{L_{\alpha,\beta}g}p{[0,\infty),x^{\alpha}dx}
    \le C \lpnorm{g}p{[0,1],x^{\alpha}dx},
    \quad g \in L^p([0,1],x^\alpha\,dx).
\end{equation*}
\end{Thm}

In what follows, we write $P^\uparrow \le Q$ with the meaning $P < Q$. In
this way, we have
\begin{equation*}
    \max\{A,B^\uparrow\} \le M \ \iff \ A \le M \text{ and } B < M.
\end{equation*}
This will keep the notation a bit shorter.

\begin{Cor}
\label{cor:seriesol}
Let $\alpha \ge -1/2$, $\beta \ge 0$, $1 < p < \infty$ and assume
\begin{equation*}
    \max\left\{ \frac{2(2\alpha+3)}{2(\alpha+\beta)+3},
        \Big(\frac{4(\alpha+1)}{2\alpha+3}\Big)^\uparrow \right\}
    \le p <\min\left\{ \frac{4(\alpha+1)}{2\alpha+1},
        \frac4{2\beta+1} \right\}.
\end{equation*}
Then, for any $g \in L^p([0,1],x^\alpha\,dx)$,
we have
\begin{equation*}
    g = \sum_{n=0}^{\infty}a_n(g)p_n^\ab,
    \quad
    a_n(g) = \int_0^1 g(x) p_n^\ab(x) x^\alpha (1-x)^\beta \, dx 
\end{equation*}
in the $L^p([0,1], x^\alpha \, dx)$-norm and
\begin{equation*}
    L_{\alpha,\beta} g = \sum_{n=0}^{\infty}b_{n}j_{n}^{\alpha+\beta},
    \quad
    b_n = 2^{-\beta} \tfrac{\Gamma(\alpha+n+1)^{1/2} (n!)^{1/2}}
        {\Gamma(\alpha+\beta+n+1)^{1/2} \Gamma(\beta+n+1)^{1/2}} a_n(g)
\end{equation*}
in the $L^p([0,\infty),x^\alpha\,dx)$-norm and almost everywhere.
\end{Cor}

For the proof of Theorem~\ref{th:acotLab} and Corollary~\ref{cor:seriesol},
see section~\ref{sec:proofLab}.

Before going on, let us write Lemma~\ref{lem:HTjn} in terms of $M_\alpha$,
$M_{\alpha,\beta}$, and $\hab$. It is clear from~(\ref{ec:HTjnab}) that
$\htra j_n^{\alpha+\beta}$ is supported on $[0,1]$, so that
\begin{equation}
     \label{ec:Majnab}
     M_{\alpha}(j_n^{\alpha+\beta}) = j_n^{\alpha+\beta}.
\end{equation}
And, taking~(\ref{ec:newja}) into account,~(\ref{ec:HTbjnab}) reads as
\begin{align}
    \label{ec:Mabjnab}
    M_{\alpha,\beta}(j_n^{\alpha+\beta})
    &= x^{-\beta} M_\alpha(t^\beta j_n^{\alpha+\beta})
       = x^{-\beta} \htra(\chi_{[0,1]}\htra (t^\beta
                          j_n^{\alpha+\beta})) \\
    &= x^{-\beta} \htra(\chi_{[0,1]} d_n p_n^\ab)
     = d_n \hab p_n^\ab
    \notag
\end{align}
with
$d_n = 2^\beta \tfrac{\Gamma(\alpha+\beta+n+1)^{1/2} \Gamma(\beta+n+1)^{1/2}}
{\Gamma(\alpha+n+1)^{1/2} (n!)^{1/2}}$.

Our main result is the following:

\begin{Thm}
\label{th:solDIE}
Let $\alpha \ge -1/2$, $0\le\beta<1$, $1<p<\infty$ and 
\begin{equation*}
    \max\left\{ \frac{2(2\alpha+3)}{2(\alpha+\beta)+3},
      \Big(\frac{4(\alpha+1)}{2\alpha+3}\Big)^\uparrow \right\}
    \le p <
    \min\left\{\frac{4(\alpha+1)}{2\alpha+4\beta+1},
                \frac{4}{2\beta+1}\right\}.
\end{equation*}
For each $g \in L^p([0,1], x^\alpha \, dx)$, $f = L_{\alpha,\beta}g$ is
a solution in $L^p([0,\infty),x^{\alpha}\,dx)$ of the dual integral
equation
\begin{equation*}
    \left\{
      \begin{array}{l}
         M_{\alpha,\beta}f = \hab g, \\[3pt]
         M_{\alpha}f = f.
      \end{array}
    \right.
\end{equation*}
\end{Thm}

\begin{proof}
Let $g \in L^p([0,1], x^\alpha \, dx)$ and $f = L_{\alpha,\beta}g$. It is
easy to see that we can apply Propositions~\ref{pro:acotMa}, \ref{pro:acotMab}
and~\ref{pro:acotHab}. Since $L_{\alpha,\beta}$ is bounded (by
Theorem~\ref{th:acotLab}), Corollary~\ref{cor:seriesol}
and~(\ref{ec:Mabjnab}) give
\begin{align*}
    M_{\alpha,\beta}f
    &= \lim_{n \to \infty} M_{\alpha,\beta}
        \Big(\sum_{k=0}^n b_kj_k^{\alpha+\beta}\Big) \\
    &= \lim_{n \to \infty} \hab
         \Big(\sum_{k=0}^n b_kd_kp_k^\ab\Big)
       = \hab g,
\end{align*}
while Corollary~\ref{cor:seriesol} and~(\ref{ec:Majnab}) yield
\begin{equation*}
    M_{\alpha}f
    = \lim_{n \to \infty}
      M_{\alpha}\Big(\sum_{k=0}^nb_kj_k^{\alpha+\beta}\Big)
    = \lim_{n \to \infty} \sum_{k=0}^n b_k j_k^{\alpha+\beta}
    = f.
   \qedhere
\end{equation*}
\end{proof}

%%%%%%%%%%%%%%%%%%%%%%%%%%%%%%%%%%%%%%%%%%%%%%%%%%%%%%%%%%%%%%%%%%%%%%%%%%%
\section{Uniqueness of the solution}
\label{sec:uniq}
%%%%%%%%%%%%%%%%%%%%%%%%%%%%%%%%%%%%%%%%%%%%%%%%%%%%%%%%%%%%%%%%%%%%%%%%%%%

Let us consider the $L^p$ subspaces
\begin{equation*}
	  B_{p,\alpha,\beta} = \overline{\spa} \{j_n^{\alpha+\beta}(x)\}_{n=0}^\infty
   \quad (\text{closure in } L^p([0,\infty),x^\alpha\,dx)),
\end{equation*}
\begin{equation*}
    E_{p,\alpha} = \left\{ f\in L^p([0,\infty),x^{\alpha}\,dx)
    : M_{\alpha}f=f \right\}.
\end{equation*}
The following results about the mean convergence of Fourier-Neumann series
were proved in~\cite{CGPV}:

\begin{Thm}
\label{th:convBp}
Let $\alpha > -1$, $\alpha+\beta > -1$, $4/3 < p < 4$, and
\begin{equation*}
    \max\left\{-\frac{\alpha+\beta+1}2, -\frac14\right\}
    < (\alpha+1) \left( \frac12 - \frac1p \right) + \frac\beta2 < \frac14.
\end{equation*}
Then, for any $f\in B_{p,\alpha,\beta}$ there exists a unique expansion
\begin{equation*}
    f = \sum_{n=0}^\infty b_n(f)j_n^{\alpha+\beta}
\end{equation*}
which holds
in the $L^p([0,\infty),x^{\alpha}\,dx)$-norm. This expansion also holds
almost everywhere.
\end{Thm}

\begin{Thm}
\label{th:Bp=Ep}
Let $\alpha \ge -1/2$, $\beta>-1/2$, $4/3<p$ with
\begin{equation*}
    \frac{-1}4 < (\alpha+1) \left( \frac12-\frac1p \right) < \frac14.
\end{equation*}
If $p<2$, assume further
\begin{equation*}
    \frac{-1}4 < (\alpha+1)\left( \frac12-\frac1p \right)-\frac{|\beta|}2.
\end{equation*}
Then $B_{p,\alpha,\beta} = E_{p,\alpha}$.
\end{Thm}

Using these results, we can prove

\begin{Thm}
Let $\alpha \ge -1/2$, $0\le\beta<1$, $1<p<\infty$ and
\begin{equation*}
    \max\left\{ \frac{2(2\alpha+3)}{2(\alpha+\beta)+3},
      \Big(\frac{4(\alpha+1)}{2\alpha-2\beta+3}\Big)^\uparrow \right\}
    \le p <
    \min\left\{\frac{4(\alpha+1)}{2\alpha+4\beta+1},
      \frac{4}{2\beta+1}\right\}.
\end{equation*}
Then, $f = L_{\alpha,\beta}g$ is the unique solution in
$L^p([0,\infty),x^{\alpha}\,dx)$ of the dual equation
\begin{equation*}
    \left\{
      \begin{array}{l}
         M_{\alpha,\beta}f = \hab g, \\[3pt]
         M_{\alpha}f = f.
      \end{array}
    \right.
\end{equation*}
\end{Thm}

\begin{proof}
It is not difficult to check that, under the hypothesis of this theorem,
we can deduce the ones of Theorems~\ref{th:convBp} and~\ref{th:Bp=Ep} and
Theorem~\ref{th:solDIE}. For instance, $\tfrac43 < p$ follows from
\begin{equation*}
    \frac43 <
    \max\left\{ \frac{2(2\alpha+3)}{2(\alpha+\beta)+3},
                 \frac{4(\alpha+1)}{2\alpha-2\beta+3} \right\}.
\end{equation*}
Indeed, if this inequality failed we would have
\begin{equation*}
   \frac{2(2\alpha+3)}{2(\alpha+\beta)+3} \leq \frac43,
   \qquad
   \frac{4(\alpha+1)}{2\alpha-2\beta+3} \leq \frac43,
\end{equation*}
which yield $2\alpha + 3 \leq 4\beta$ and $\alpha + 2\beta \leq 0$,
respectively. Then $\alpha \leq -3/4$, which contradicts $\alpha \geq
-1/2$.

According to Theorem~\ref{th:solDIE}, $L_{\alpha,\beta}g$ is a solution of the
dual equation. Let us see that it is unique. Let $f$ be a solution, that
is, $f \in  L^p([0,\infty),x^{\alpha}\,dx)$, $M_\alpha f = f$ and
$M_{\alpha,\beta}f = \hab g$. In particular, $f \in E_{p,\alpha}$. By
Theorems~\ref{th:Bp=Ep} and~\ref{th:convBp}, we can expand~$f$ as a
Fourier-Neumann series
\begin{equation*}
    f = \sum_{n=0}^{\infty} b_n(f) j_n^{\alpha+\beta}
\end{equation*}
which converges in the $L^p([0,\infty),x^\alpha\,dx)$-norm. Proving that each
$b_n(f)$ is uniquely determined will suffice.

{}From~(\ref{ec:Mabjnab}) and the fact that $M_{\alpha,\beta}$ is a
continuous (bounded) operator in $L^p$, we get
\begin{equation*}
   \hab g = M_{\alpha,\beta}f
    = \sum_{n=0}^{\infty} b_n(f)d_n \hab p_n^\ab.
\end{equation*}
Our assumptions on $\alpha$, $\beta$, and $p$, togehter with the estimates
\begin{equation*}
    J_\alpha(x) = O(x^\alpha), \quad x\to0^+,
    \qquad\text{and}\qquad
    J_\alpha(x) = O(x^{-1/2}), \quad x\to\infty,
\end{equation*}
yield $x^\beta j_k^{\alpha+\beta} \in L^{p'}([0,\infty),x^\alpha\,dx)$, where
$1/p+1/p'=1$. Hence, the map $h \mapsto \int_0^\infty x^\beta
j_k^{\alpha+\beta}(x)h(x)x^\alpha\,dx$ is a continuous operator from
$L^p([0,\infty),x^\alpha\,dx)$ into $\mathbb{R}$. Then,
\begin{multline*}
    \int_0^\infty x^\beta j_k^{\alpha+\beta}(x) \hab(g,x) x^\alpha\,dx
    \\
    = \sum_{n=0}^\infty b_n(f) d_n
    \int_0^\infty x^\beta j_k^{\alpha+\beta}(x) \hab(p_n^\ab,x) x^\alpha\,dx.
\end{multline*}
Now, recall the multiplication formula for the Hankel transform, which is
valid for $h_1, h_2 \in L^2([0,\infty),x^\alpha\,dx)$:
\begin{equation*}
    \int_0^\infty h_1(x) \htra(h_2,x) x^\alpha\,dx
    = \int_0^\infty \htra(h_1,x) h_2(x)x^\alpha\,dx.
\end{equation*}
Indeed, in $S^+$ this follows from Fubini's theorem, then it extends to the
whole $L^2([0,\infty),x^\alpha\,dx)$ by continuity.

Thus, the definition of $\hab$, together with (\ref{ec:HTjnab}) and the
orthogonality of Jacobi polynomials, yield
\begin{gather*}
    \int_0^\infty x^\beta j_k^{\alpha+\beta}(x) \hab(p_n^\ab,x) x^\alpha\,dx
    = \int_0^\infty j_k^{\alpha+\beta}(x) \htra(\chi_{[0,1]}p_n^\ab,x)
          x^\alpha\,dx \\
    = r_k \int_0^1 (1-x)^\beta p_k^\ab(x) p_n^\ab(x) x^\alpha\,dx
    = r_k \delta_{kn}
\end{gather*}
with a constant $r_k\ne0$ (actually, $r_k = 1/d_k$, where $d_k$
comes from~(\ref{ec:Mabjnab}), as before). Therefore,
\begin{equation*}
    \int_0^\infty x^\beta j_n^{\alpha+\beta}(x) \hab g(x) x^\alpha\,dx
   = b_n(f).
   \qedhere
\end{equation*}
\end{proof}

%%%%%%%%%%%%%%%%%%%%%%%%%%%%%%%%%%%%%%%%%%%%%%%%%%%%%%%%%%%%%%%%%%%%%%%%%%%
\section{Proof of Lemma~\ref{lem:HTjn}}
\label{sec:proofHTjn}
%%%%%%%%%%%%%%%%%%%%%%%%%%%%%%%%%%%%%%%%%%%%%%%%%%%%%%%%%%%%%%%%%%%%%%%%%%%

As we already mentioned, the first part of Lemma~\ref{lem:HTjn} was proved
in~\cite{CGPV}, so we only prove now the second part. Let $\hf$ denote, as
usual, the hypergeometric function. We use the formula
\begin{multline*}
    \int_0^\infty t^{-\lambda} J_\mu(at) J_\nu(bt)\,dt
    \\
    = \tfrac{2^{-\lambda} \Gamma(\frac12(\mu+\nu-\lambda+1))b^\nu}
       {a^{\nu-\lambda+1} \Gamma(\nu+1)\Gamma(\frac12(\lambda+\mu-\nu+1))}
     \hf\left( \tfrac{\mu+\nu-\lambda+1}2,
       \tfrac{\nu-\lambda-\mu+1}2; \nu+1; \tfrac{b^2}{a^2} \right),
     \ 0<b<a,
\end{multline*}
valid when $\mu+\nu-\lambda>-1$ and $\lambda>-1$ (see~\cite[8.11 (9),
p.~48]{E-TIT} or~\cite[13.4 (2), p. 401]{Wat}). Take $a=1$ and $x=b^2$, with
parameters $\lambda=-\beta$, $\mu=\alpha+\beta+2n+1$ and $\nu=\alpha$. Then,
for $\beta<1$, $\alpha+\beta>-1$, and $0 < x < 1$ we get, after a change of
variable,
\begin{multline*}
   \frac{x^{-\alpha/2}}{2} \int_0^\infty y^\beta j_n^{\alpha+\beta}(y)
   J_\alpha(\sqrt{xy}) y^{\alpha/2} \, dy
   \\
    = \tfrac{2^\beta \sqrt{\alpha+\beta+2n+1}\,\Gamma(\alpha+\beta+n+1)}
      {\Gamma(\alpha+1) \Gamma(n+1)}
      \hf(\alpha+\beta+n+1,-n;\alpha+1;x).
\end{multline*}
Taking into account that
\begin{equation*}
    P_n^{(\alpha,\beta)}(x) =
      \tfrac{\Gamma(\alpha+n+1)}{\Gamma(\alpha+1)\Gamma(n+1)}
      \hf(\alpha+\beta+n+1, -n ; \alpha+1; \tfrac{1-x}{2})
\end{equation*}
whenever $\alpha, \beta > -1$ and $-1 < x < 1$, it follows
\begin{multline*}
   \frac{x^{-\alpha/2}}{2} \int_0^\infty y^\beta j_n^{\alpha+\beta}(y)
   J_\alpha(\sqrt{xy}) y^{\alpha/2} \, dy
   \\
    = \tfrac{2^\beta \sqrt{\alpha+\beta+2n+1}\,\Gamma(\alpha+\beta+n+1)}
      {\Gamma(\alpha+n+1)}
      P_n^{(\alpha,\beta)}(1-2x),
\end{multline*}
if $x \in (0,1)$. We have not finished yet, because this integral must be
understood as an improper Riemann integral, not a Lebesgue integral. In other
words, this means
\begin{equation*}
   \lim_{R\to\infty} \htra(t^\beta j_n^{\alpha+\beta} \chi_{[0,R]},x)
    = \tfrac{2^\beta \sqrt{\alpha+\beta+2n+1}\,\Gamma(\alpha+\beta+n+1)}
      {\Gamma(\alpha+n+1)}
      P_n^{(\alpha,\beta)}(1-2x),
\end{equation*}
pointwisely on $(0,1)$. Now, $\htra$ is a bounded operator, so that
\begin{equation*}
   \lim_{R\to\infty} \htra(t^\beta j_n^{\alpha+\beta} \chi_{[0,R]})
   = \htra(t^\beta j_n^{\alpha+\beta})
\end{equation*}
in the $L^p$-norm. It follows that
\begin{equation*}
   \htra(t^\beta j_n^{\alpha+\beta},x)
    = \tfrac{2^\beta \sqrt{\alpha+\beta+2n+1}\,\Gamma(\alpha+\beta+n+1)}
      {\Gamma(\alpha+n+1)}
      P_n^{(\alpha,\beta)}(1-2x),
\end{equation*}
almost everywhere on $[0,1]$.

%%%%%%%%%%%%%%%%%%%%%%%%%%%%%%%%%%%%%%%%%%%%%%%%%%%%%%%%%%%%%%%%%%%%%%%%%%%
\section{Proof of Theorem~\ref{th:acotLab} and Corollary~\ref{cor:seriesol}}
\label{sec:proofLab}
%%%%%%%%%%%%%%%%%%%%%%%%%%%%%%%%%%%%%%%%%%%%%%%%%%%%%%%%%%%%%%%%%%%%%%%%%%%

\begin{Lem}
\label{lem:acotopG}
Let $\alpha \geq {-1/2}$, $\beta\ge0$, and $1 < p < \infty$, with
$\tfrac{2(2\alpha+3)}{2(\alpha+\beta)+3} \leq p$.
Then
\begin{equation*}
  \|\htr_{\alpha+\beta}(\chi_{[0,1]} h)\|_{L^p([0,\infty),x^\alpha\,dx)}
   \leq C \| h \|_{L^p([0,1],x^\alpha\,dx)},
   \quad h \in L^p([0,1],x^\alpha\,dx).    
\end{equation*}
\end{Lem}

\begin{proof}
Take $\nu = \alpha + \beta + \tfrac32 - 2\tfrac{\alpha+1}p$. Then
\begin{equation*}
  \|\htr_{\alpha+\beta}(\chi_{[0,1]} h)\|_{L^p([0,\infty),x^\alpha\,dx)}
   = \| x^{-\nu/2 + (\alpha+\beta)/2 + 3/4}
   \htr_{\alpha+\beta}(\chi_{[0,1]} h)\|_{L^p([0,\infty),\tfrac{dx}x)}.
\end{equation*}
With our assumptions on $\alpha$, $\beta$, and $p$ we can apply
Theorem~\ref{th:Rooney} and get
\begin{align*}
  \|\htr_{\alpha+\beta}(\chi_{[0,1]} h)\|_{L^p([0,\infty),x^\alpha\,dx)}
  &\leq C \| x^{\nu/2 + (\alpha+\beta)/2 + 1/4}
   \chi_{[0,1]} h\|_{L^p([0,\infty),\tfrac{dx}x)}
   \\
  &= C 
  \| x^{\alpha+\beta+1-2(\alpha+1)/p} h\|_{L^p([0,1],x^\alpha \, dx)}.
\end{align*}
The easy observation that $\alpha+\beta+1-2\tfrac{\alpha+1}p \geq 0$ finishes
the proof.
\end{proof}

\begin{proof}[Proof of Theorem~\ref{th:acotLab}]
Lemma~\ref{lem:acotopG} proves Theorem~\ref{th:acotLab} in the case $\beta =
0$, since $L_{\alpha,0}g = \htra(\chi_{[0,1]} g)$. Now, observe that
$L_{\alpha,\beta} g = 2^{-\beta} \htr_{\alpha+\beta} (\chi_{[0,1]}
I_{\alpha,\beta}g)$ if $\beta > 0$, where
\begin{equation*}
    I_{\alpha,\beta}(g,x)
    = \frac{x^{-(\alpha+\beta)}}{\Gamma(\beta)}
      \int_{0}^x g(t)(x-t)^{\beta-1}t^{\alpha}\,dt,
    \quad 0 < x < 1
\end{equation*}
is the Erd\'elyi-Kober operator. It is well known that this operator is
bounded in $L^p([0,1], x^\alpha \, dx)$ if $\alpha > -1$, $\beta > 0$, and $1
< p < \infty$. Indeed, after a change of variable we obtain
\begin{equation*}
    I_{\alpha,\beta}(g,x)
    = \frac1{\Gamma(\beta)}\int_0^1(1-z)^{\beta-1}z^{\alpha}g(xz)\,dz
\end{equation*}
and, by Minkowski's integral inequality,
\begin{align*}
    \lpnorm{I_{\alpha,\beta}g}p{[0,1],x^{\alpha}dx}
    &\le \frac1{\Gamma(\beta)} \int_0^1(1-z)^{\beta-1}z^{\alpha} 
       \lpnorm{g(xz)}p{[0,1],x^{\alpha}dx}\,dz \\
    &\le \lpnorm{g}p{[0,1],x^{\alpha}dx} \frac1{\Gamma(\beta)}
       \int_0^1(1-z)^{\beta-1}z^{\alpha-{(\alpha+1)}/p}\,dz \\
    &= C \lpnorm{g}p{[0,1],x^{\alpha}dx},
\end{align*}
where we have used that $\lpnorm{g(xz)}p{[0,1],x^{\alpha}dx} \le
z^{-(\alpha+1)/p}\lpnorm{g}p{[0,1],x^{\alpha}dx}$ and
\begin{equation*}
    \int_0^1(1-z)^{\beta-1}z^{\alpha-{(\alpha+1)}/p}\,dz
    = B(\beta,(\alpha+1)(1-1/p)) < \infty.
\end{equation*}
Thus, $I_{\alpha,\beta}$ is bounded and again Lemma~\ref{lem:acotopG} proves
the theorem.
\end{proof}

\begin{proof}[Proof of~Corollary~\ref{cor:seriesol}]
Let $g \in L^p([0,1], x^\alpha \, dx)$. Under our assumptions on $\alpha$,
$\beta$, and $p$, it is easy to check that we can apply
Corollary~\ref{cor:convja}. Therefore,
\begin{equation*}
    g = \sum_{n=0}^{\infty}a_n(g)p_n^\ab,
    \quad
    a_n(g) = \int_0^1 g(x) p_n^\ab(x) x^\alpha (1-x)^\beta \, dx,
\end{equation*}
in the $L^p([0,1], x^\alpha \, dx)$-norm.
By Theorem~\ref{th:acotLab}, $L_{\alpha,\beta}$ is a continuous (i.e. bounded)
operator from $L^p([0,1], x^\alpha \, dx)$ into $L^p([0,\infty), x^\alpha \,
dx)$. Then,
\begin{equation*}
   L_{\alpha,\beta} g = \sum_{n=0}^{\infty}a_n(g) L_{\alpha,\beta} p_n^\ab,
\end{equation*}
where the convergence holds in the $L^p$-norm. Now, consider the following
formula (see~\cite[13.1 (43), p.~191]{E-TIT}):
\begin{equation*}
    I_{\alpha,\beta}(P_n^\ab(1-2t),x)
    = \tfrac{\Gamma(\alpha+n+1)}{\Gamma(\alpha+\beta+n+1)}
      P_n^{(\alpha+\beta,0)}(1-2x).
\end{equation*}
Lemma~\ref{lem:HTjn} (with parameters $\alpha+\beta$ and $0$ instead of
$\alpha$ and $\beta$, respectively) gives
$\htr_{\alpha+\beta}(j_n^{\alpha+\beta},x) = \sqrt{\alpha+\beta+2n+1}
P_n^{(\alpha+\beta,0)}(1-2x) \chi_{[0,1]}(x)$, so that
\begin{equation}
   \label{ec:Ha+bPa+b0}
   \htr_{\alpha+\beta}( \chi_{[0,1]} P_n^{(\alpha+\beta,0)}(1-2t))
   = (\alpha+\beta+2n+1)^{-1/2} j_n^{\alpha+\beta}
\end{equation}
(since $\htr_{\alpha+\beta}^2 = \Id$ in $L^2$). Thus,
\begin{equation}
    \label{ec:Labja}
    L_{\alpha,\beta}(P_n^{(\alpha,\beta)}(1-2t))
    = 2^{-\beta}
  \frac{\Gamma(\alpha+n+1)}{\sqrt{\alpha+\beta+2n+1} \Gamma(\alpha+\beta+n+1)}
   j_n^{\alpha+\beta}
\end{equation}
if $\beta > 0$. In the case $\beta = 0$, (\ref{ec:Ha+bPa+b0}) and
$L_{\alpha,0}g = \htra(\chi_{[0,1]} g)$ give (\ref{ec:Labja}), as well.
In terms of the normalised polynomials $p_n^{(\alpha,\beta)}$, this means
\begin{equation*}
    L_{\alpha,\beta} p_n^\ab =
      2^{-\beta} \tfrac{\Gamma(\alpha+n+1)^{1/2} (n!)^{1/2}}
        {\Gamma(\alpha+\beta+n+1)^{1/2} \Gamma(\beta+n+1)^{1/2}}
      j_n^{\alpha+\beta},
\end{equation*}
so that
\begin{equation*}
    L_{\alpha,\beta} g =
    \sum_{n=0}^\infty a_n(g)
      2^{-\beta} \tfrac{\Gamma(\alpha+n+1)^{1/2} (n!)^{1/2}}
        {\Gamma(\alpha+\beta+n+1)^{1/2} \Gamma(\beta+n+1)^{1/2}}
      j_n^{\alpha+\beta}
\end{equation*}
in the $L^p$-norm and, by Theorem~\ref{th:convBp}, almost everywhere.
\end{proof}

%%%%%%%%%%%%%%%%%%%%%%%%%%%%%%%%%%%%%%%%%%%%%%%%%%%%%%%%%%%%%%%%%%%%%%%%%%%
%References
%%%%%%%%%%%%%%%%%%%%%%%%%%%%%%%%%%%%%%%%%%%%%%%%%%%%%%%%%%%%%%%%%%%%%%%%%%%

\begin{thebibliography}{99}

\bibitem{CGPV} \'O. Ciaurri, J. J. Guadalupe, M. P\'erez and J. L. Varona,
\textit{Mean and almost everywhere convergence of Fourier-Neumann series},
{J. Math. Anal. Appl.}
{236} (1999), 125--147.
 
\bibitem{Dur} A. J. Dur\'an,
\textit{On Hankel transform},
{Proc. Amer. Math. Soc.}
{110} (1990), 417--424.

\bibitem{E-HTF} A. Erd\'elyi, W. Magnus, F. Oberhettinger and F. Tricomi,
\textit{Higher Transcendental Functions,} Vol.~II,
McGraw-Hill, New York, 1953.

\bibitem{E-TIT} A. Erd\'elyi, W. Magnus, F. Oberhettinger and F. Tricomi,
\textit{Tables of Integral Transforms,} Vol.~II,
McGraw-Hill, New York, 1954.

\bibitem{Her} C. S. Herz,
\textit{On the mean inversion of Fourier and Hankel transforms},
{Proc. Nat. Acad. Sci. U.S.A.}
{40} (1954), 996--999.

\bibitem{H-R} P. Heywood and P. G. Rooney,
\textit{A weighted norm inequality for the Hankel transformation},
{Proc. Roy. Soc. Edinburgh Sect. A}
{99} (1984), 45--50.

\bibitem{Hoc} H. Hochstadt,
\textit{Integral equations,}
Wiley, New York, 1973.

\bibitem{Muc} B. Muckenhoupt,
\textit{Mean convergence of Jacobi series},
{Proc. Amer. Math. Soc.}
{23} (1969), 306--310.

\bibitem{Roo} P. G. Rooney,
\textit{A technique for studying the boundedness and extendability of certain
types of operators},
{Canad. J. Math.}
{25} (1973), 1090--1102.

\bibitem{Sne} I. N. Sneddon,
\textit{Fourier Transforms,}
McGraw-Hill, New York, 1951.
Republication: Dover, New York, 1995.

\bibitem{Ste} K. Stempak,
\textit{Transplanting maximal inequalities between Laguerre and Hankel
multipliers},
{Monatsh. Math.}
{122} (1996), 187--197.

\bibitem{St-Tr} K. Stempak and W. Trebels,
\textit{Hankel multipliers and transplantation operators},
{Studia Math.}
{126} (1997), 51--66.

\bibitem{Sze} G. Szeg\H{o},
\textit{Orthogonal Polynomials} (4th edition),
Amer. Math. Soc. Colloq. Publ.
{23}, Amer. Math. Soc.,
Providence, RI, 1975.

\bibitem{Tit} E. C. Titchmarsh,
\textit{Introduction to the Theory of Fourier Integrals,}
Oxford, New York, 1937.

\bibitem{Tra1} C. J. Tranter,
\textit{Integral Transforms in Mathematical Physics} (3th edition),
Methuen, London, 1966.

\bibitem{Tra2} C. J. Tranter,
\textit{Bessel Functions with Some Physical Applications,}
English Univ. Press, London, 1968.

\bibitem{Var} J. L. Varona,
\textit{Fourier series of functions whose Hankel transform is supported on
$[0,1]$},
{Constr. Approx.}
{10} (1994), 65--75.

\bibitem{Wat} G. N. Watson,
\textit{A Treatise on the Theory of Bessel Functions,}
Cambridge Univ. Press, 1966.

\end{thebibliography}

\end{document}